\section{第\ref{sec:gaba}章　\nt{GLUT}と\nt{GABA}の協働}

この章の論理的・力学的な主張は、線形解析とオームの法則から章の中で導いたものである。実験が示すのは、それらが依拠する回路、つまり遅延つきの禁止、興奮に続く直接の抑制、抑制による安定化、\nt{GLUT}--\nt{GABA}ループのリズムが、脳で実際に働いていることである。

{\footnotesize
\setlength{\tabcolsep}{3pt}\renewcommand{\arraystretch}{1.15}
\begin{longtable}{>{\raggedright}p{0.5cm}>{\raggedright}p{4.0cm}>{\raggedright}p{5.6cm}>{\raggedright}p{2.3cm}>{\raggedright\arraybackslash}p{1.7cm}}
\toprule
No. & 主張 & 実験と測定されたもの & 出典 & 状態 \\
\midrule\endfirsthead
\toprule
No. & 主張 & 実験と測定されたもの & 出典 & 状態 \\
\midrule\endhead
1 & \nt{GABA}ニューロンは皮質のニューロンの5分の1から4分の1 & サル皮質の$10$領域でGABAを免疫染色：$8$領域でニューロンの約$25\,\%$、一次視覚野で$20\,\%$未満。皮質の抑制性ニューロンの多様性についての総説 & \cite{Hendry1987,Markram2004} & 測定済み \\
2 & 抑制が何をするかは着地点で大きく決まる。樹状突起、細胞体、スパイクの発生点、他の配線の終末 & 総説：皮質の\nt{GABA}ニューロンのクラスは受け手上の標的で区別される。錐体ニューロンの細胞体と樹状突起からの同時記録：直接の抑制は樹状突起より細胞体ではるかに強い。脊髄で他の配線の終末への抑制は、1本の入力線の神経伝達物質放出を減らす（測定の総説） & \cite{Markram2004,Pouille2001,Rudomin1999} & 測定済み \\
3 & 仲介役を通じた符号の反転には遅延1段、約1ミリ秒がかかる。興奮に続いて届く抑制は同時性の窓を狭める & ラット海馬スライスの錐体ニューロン：仲介役1個を介した抑制は直接の興奮に短い遅延で続き、二つの興奮性入力が閾値まで加算される窓は$2\ms$未満。この抑制が弱い樹状突起では窓が広い & \cite{Pouille2001} & 測定済み \\
4 & 禁止$a\wedge\bar b$~\eqref{eq:veto}はORと定数1とともに関数的に完全である（命題~\ref{prop:gabacomplete}） & 章内の証明。古典的な結果：抑制性入力をもつ閾値素子のネットワークは任意の論理式を実現する。モデルについての主張であり、実験はない & \cite{McCullochPitts1943} & 理論 \\
5 & 遅延つきの禁止は、最適速度$v^*=\Delta x/d$~\eqref{eq:vstar}の運動方向検出器になる & ウサギの網膜：方向選択性は、「ヌル」方向の運動で興奮を打ち消す抑制によって作られる。選択性をもつ細胞の興奮性入力と抑制性入力を別々に記録すると、スターバーストアマクリン細胞（\nt{GABA}）が非対称に、つまり「ヌル」側に向いた突起だけで抑制していることがわかった。$v^*$の式は章の導出 & \cite{Barlow1965,Fried2002} & 測定済み \\
6 & シャント抑制は電圧から引くのではなく割る~\eqref{eq:shunt}（図~\ref{fig:shunt}） & 塩化物イオンの平衡電位が静止電位付近にあるときの、三つのコンダクタンスからなる分圧器のオームの法則。スパイクを出さないニューロンについて & 章内の導出 & 理論 \\
7 & スパイクの発火率に対してシャントは引き算として作用し、割り算は抑制と釣り合った背景ノイズの組み合わせから得られる & モデル：発火しているニューロンではシャントを通る電流はほぼ一定で、発火率への効果は減算的。実験：ラット体性感覚野の錐体ニューロン（スライス）にダイナミッククランプで背景の興奮性・抑制性コンダクタンスを注入した。両者を釣り合わせて同時に増やすと、発火率をあまりずらさずに応答の傾きが割り算される。総説：全体の活動による割り算は脳の多くの部位で見られる & \cite{HoltKoch1997,Chance2002,Carandini2012} & 測定済み \\
8 & $\beta>1$では相互抑制がただ一つの勝者を選ぶ（命題~\ref{prop:wta}、~\eqref{eq:wta}） & 固有値$-(1\pm\beta)/\tau$は章の導出。$\beta=0.5$での発火率$0.73$と$0.53$は固定点$r_{1,2}=(I_{1,2}-\beta I_{2,1})/(1-\beta^2)$。\texttt{models/}にスクリプトはない。章では直接の実験は挙げていない & 章内の導出 & 理論 \\
9 & \nt{GABA}を介した信号によるメモリのリセット。相互抑制する二つの群はラッチである & 図~\ref{fig:bistable}と命題~\ref{prop:wta}から従う。実験はない & 章内の導出 & 理論 \\
10 & \nt{GLUT}--\nt{GABA}ループの平衡は、抑制が十分に強く十分に速ければ安定である（命題~\ref{prop:ei}） & 二集団モデル~\eqref{eq:ei}。条件(i)と(ii)はその行列の行列式とトレースであり、章で導いた & \cite{WilsonCowan1972} & 理論 \\
11 & $w_{EE}=1.5$、$w_{EI}w_{IE}=2$のとき、速い抑制（$\tau_I=2\ms$）は$E^*=0.67h$を保ち、遅い抑制（$30\ms$）は振動を起こす（図~\ref{fig:ei}） & 式~\eqref{eq:ei}の数値解、スクリプト\texttt{figures/make\_ei.py} & 計算 & モデル \\
12 & 皮質は抑制安定化の状態で動作する。その特徴は、\nt{GABA}ニューロンを後押しするとその発火率が下がること~\eqref{eq:isn} & 理論が逆説的な応答を予言した。実験：ネコ一次視覚野の細胞内記録で、周囲の刺激で応答が抑えられるとき、抑制性コンダクタンスは増えずに興奮性コンダクタンスと一緒に減る。マウスの視覚野、体性感覚野、運動野でPVニューロンを光遺伝学的に興奮させると、刺激がなくても、また軽い麻酔下でも、抑制性ニューロンの発火率が下がる。図~\ref{fig:ei}の値での係数$-1/3$は章の導出 & \cite{Tsodyks1997isn,Ozeki2009,Sanzeni2020} & 測定済み \\
13 & 「\nt{GLUT}が\nt{GABA}を興奮させ、\nt{GABA}が\nt{GLUT}を抑制する」ループは周期$12$--$50\ms$の速いリズムを生む & マウス体性感覚野の速い\nt{GABA}ニューロンを$8$--$200$\,Hzで光遺伝学的に刺激すると、ガンマリズム（$20$--$80$\,Hz、周期$12$--$50\ms$）が選択的に強まる。\nt{GLUT}ニューロンの刺激はより遅いリズムだけを強める & \cite{Cardin2009,Buzsaki2012} & 測定済み \\
\bottomrule
\end{longtable}}
